\section{Research Gap and Positioning of Motivo}\label{sec:motivo-position}

The systems discussed in this chapter illustrate a wide spectrum of approaches to music programming and representation.
Live-coding languages such as Sonic Pi, Tidal, Strudel, and HMusic prioritize interactive performance and pattern
transformation, typically streaming events to a synthesis backend rather than producing a self-contained offline
artifact~\cite{Aaron_2016,McLean_Wiggins_2010,Roos_McLean_2023,Du_Bois_Ribeiro_2019}.
Real-time audio languages including ChucK and SuperCollider provide precise control over synthesis and timing, but at a
level of abstraction oriented toward signal graphs and concurrent processes rather than hierarchical song
structure~\cite{Wang_Cook_2003,McCartney_2002}.
Interchange formats such as MIDI and ABC notation encode musical information for playback or typesetting, yet neither
provides the executable abstractions---typed patterns, loops, and conditional generation---that Motivo offers at the
source level~\cite{Moog_1986,Walshaw_2011,Moore_1988}.

Table~\ref{tab:related-work-comparison} summarizes the main systems and highlights the design trade-offs that motivate
Motivo.

\begin{table}[htbp]
  \centering
  \footnotesize
  \renewcommand{\arraystretch}{1.5}
  \begin{tabular}{p{2.2cm}p{2.5cm}p{4cm}p{5cm}}
    \toprule
    \textbf{System / Language} & \textbf{Class} & \textbf{Temporal / structural model} & \textbf{Main gap
    relative to this thesis} \\
    \midrule
    Sonic Pi & Live coding & Sequential execution with notes, sleeps, and simple patterns &
    Live synthesis target; limited structural abstraction for offline MIDI compilation \\
    Tidal / Strudel & Live coding & Time-varying patterns with functional composition &
    Event streaming to external backends; no deterministic offline MIDI as primary artifact \\
    HMusic & Live coding & Grid-like patterns and parallel tracks & Compiles to Sonic Pi; remains performance-coupled \\
    ChucK & Real-time synthesis & Explicit time advancement and synchronized concurrent shreds &
    Low-level time management; weak high-level compositional structure \\
    Super\-Collider & Real-time synthesis & UGens, patterns, and streams over OSC &
    Synthesis-oriented; composition and execution not separated \\
    MIDI & Interchange format & Discrete message stream of note and control events &
    Strong backend format, too low-level for structured source-level composition \\
    ABC notation & Notation format & Declarative textual score with implicit ordering &
    Not an executable language; no patterns or control flow \\
    \bottomrule
  \end{tabular}
  \caption{Comparison of the main systems and languages discussed in this
  chapter.}\label{tab:related-work-comparison}
\end{table}

From this analysis, three gaps relative to the goals of this thesis emerge.
First, most musical programming languages target live performance or synthesis, leaving offline deterministic MIDI
compilation underexplored as a primary design goal.
Second, there is a gap between low-level audio programming languages and high-level compositional tools: few systems
provide both expressive procedural structure and a zero-runtime compilation model.
Third, existing approaches often conflate composition with execution, making it difficult to separate musical intent
from performance-time scheduling.

Motivo addresses these gaps by providing a typed, imperative source language for offline symbolic composition;
a compiler that validates structure and expands it deterministically into a flat event timeline; and a Standard MIDI
File as the portable output artifact.
Patterns and loops compress repetitive structure in the source program while the lowering stage materializes the full
event stream---the compression argument introduced in Section~\ref{sec:algorithmic-composition}.

The next chapter moves from the surrounding ecosystem to the design of Motivo itself. It explains the
language model, the compiler architecture, and the lowering process that turns recursive musical structure
into MIDI events.
